\centerline{\bf \Huge Тренировочная олимпиада}
\centerline{\bf \Huge 10 класс}

\begin{flushright}
        \scriptsize{Арбуз арбуз пока}
\end{flushright}

\problem{1}
При $x, y, z \in[0,1]$ докажите неравенство $x y+y z+z x \leqslant 2 x y z+1$.

\problem{2}
Решите функциональное уравнение $f\left(x^2-y\right)=f\left(x^3+2 y\right)+f\left(x^4\right)$.

\problem{3}
В ЭлМШ учится 40 красоток и 10 красавчиков. Каждый красавчик влюблен ровно в тех красоток, которые умнее его, либо ровно в тех, которые богаче его. Докажите, что у каких-то двух красоток одни и те же поклонники. (Про любых двух людей можно сказать, кто из них умнее. Если человек A умнее человека B и человек B умнее человека C, то верно, что A умнее C. Такие же свойства выполнены для богатства.)

\problem{4}
Во вписанном четырехугольнике $ABCD$ биссектрисы углов $A, B, C$ и $D$ пересекают стороны четырехугольника в точках $K, L, M$ и $N$ соответственно. Докажите, что $KLMN$ вписанный.

\problem{5}
Даяна выписала на доску $57$ различных натуральных чисел, не превосходящих
$2024$. Расул Владимирович заметил, что среди любых трёх из этих чисел можно выбрать два, которые в произведении дадут точный квадрат. Докажите, что среди чисел, выписанных Даяной на доску, встретится точный квадрат.